\clearpage\thispagestyle{empty}
\pdfbookmark[0]{Quellen zum zweiten Kometen von 1805}{comet1805-2-sources}
\begin{center}\textbf{Quellen zum zweiten Kometen von 1805}\end{center}
\noindent Diese KI-Übertragung in heutiges Deutsch folgt dem Auszug aus Franz Xaver von Zachs Bericht in \href{https://gdz.sub.uni-goettingen.de/id/PPN236007467}{Gauß’ \textit{Werken}, Band~VI} (Göttingen, 1874), S.~266. Der Bericht erschien als Beitrag~VI in der \textit{Monatlichen Correspondenz}, Band~XIII, Januar 1806, \href{https://www.digitale-sammlungen.de/en/view/bsb10538605?page=91}{S.~83--91}. Der Auszug entspricht der Passage, die auf \href{https://www.digitale-sammlungen.de/en/view/bsb10538605?page=96}{S.~88} mit Carl Friedrich Gauß’ Beobachtungen vom 8. Dezember beginnt. Auf \href{https://www.digitale-sammlungen.de/en/view/bsb10538605?page=97}{S.~89} folgen seine parabolischen Elemente und Formeln, auf \href{https://www.digitale-sammlungen.de/en/view/bsb10538605?page=98}{S.~90} seine Abstände und Lichtstärken. Die genannte Methode verweist auf Gauß’ Aufsatz vom Mai 1804 auf \hyperlink{native-99}{S.~94} und auf die zugehörigen Formeln in seinem Aufsatz zum Tierkreis der Juno auf \hyperlink{native-125}{S.~120} dieses Bandes. Übertragung und Satz wurden mit KI erstellt. Die Zeitschriftendigitalisate stammen von der Bayerischen Staatsbibliothek.
\clearpage
